\documentclass[11pt,a4paper]{article}
\usepackage[utf8]{inputenc}
\usepackage[T1]{fontenc}
\usepackage[provide=*,french]{babel}
\usepackage{amsmath,amssymb,mathtools,amsthm}
\usepackage[most]{tcolorbox}
\usepackage{geometry,enumitem,hyperref}
\usepackage[expansion=false]{microtype}
\geometry{margin=2.5cm}
\hypersetup{colorlinks=true,linkcolor=blue!55!black}
\newcommand{\e}{\mathrm e}
\newcommand{\N}{\mathbb N}
\newcommand{\eps}{\varepsilon}
\newtcolorbox{theorembox}{breakable,colback=blue!3,colframe=blue!60!black,title=Théorème}
\newtcolorbox{lemmabox}[1][]{breakable,colback=green!3,colframe=green!50!black,title=Lemme,#1}
\newtcolorbox{propositionbox}[1][]{breakable,colback=orange!4,colframe=orange!70!black,title=Proposition,#1}
\newtcolorbox{remarkbox}{breakable,colback=violet!3,colframe=violet!55!black,title=Remarque}
\newtcolorbox{proofbox}[1][]{breakable,colback=gray!2,colframe=gray!65!black,title=Démonstration,#1}
\title{Exercices 25 et 26 de Dieudonné\\\large Plusieurs démonstrations rigoureuses}
\author{}
\date{}

\begin{document}
\maketitle
\tableofcontents

\section{Exercice 25 : somme de grandes puissances}

\begin{theorembox}
Soit $\alpha>0$. Lorsque $n\to\infty$,
\[
 \sum_{m=1}^{n}m^{\alpha n}\sim
 \frac{n^{\alpha n}}{1-\e^{-\alpha}}.
\]
\end{theorembox}

Posons
\[
 A_n:=n^{-\alpha n}\sum_{m=1}^{n}m^{\alpha n}
     =\sum_{k=0}^{n-1}\left(1-\frac{k}{n}\right)^{\alpha n},
\]
où l'on a effectué le changement d'indice $k=n-m$. Il suffit de prouver
$A_n\to(1-\e^{-\alpha})^{-1}$.

\subsection{Première preuve : découpage à $\lfloor n^\lambda\rfloor$}

Cette preuve suit l'indication et n'utilise aucune convergence dominée discrète.

\begin{lemmabox}[title=Approximation locale quantitative]
Pour $0\leq k\leq n/2$,
\[
0\leq \e^{-\alpha k}-\left(1-\frac{k}{n}\right)^{\alpha n}
\leq \frac{\alpha k^2}{n}\e^{-\alpha k}.
\]
\end{lemmabox}

\begin{proofbox}
Pour $0\leq x\leq 1/2$,
\[
0\leq-\log(1-x)-x=\sum_{j=2}^{\infty}\frac{x^j}{j}
\leq\frac{x^2}{2(1-x)}\leq x^2.
\]
Avec $x=k/n$, posons
\[
u_{n,k}:=\alpha n\left(-\log(1-k/n)-k/n\right).
\]
Alors $0\leq u_{n,k}\leq\alpha k^2/n$ et
\[
\left(1-\frac{k}{n}\right)^{\alpha n}=\e^{-\alpha k}\e^{-u_{n,k}}.
\]
L'inégalité $0\leq1-\e^{-u}\leq u$ pour $u\geq0$ donne le résultat.
\end{proofbox}

\begin{proofbox}[title=Preuve par découpage]
Fixons $0<\lambda<1$ et $K_n=\lfloor n^\lambda\rfloor$. Pour $n$ assez
grand, $K_n\leq n/2$. D'une part, le lemme donne
\[
\sum_{k=0}^{K_n}\left|\left(1-\frac{k}{n}\right)^{\alpha n}
-\e^{-\alpha k}\right|
\leq\frac{\alpha}{n}\sum_{k=0}^{\infty}k^2\e^{-\alpha k}=O(n^{-1}).
\]
D'autre part, $\log(1-x)\leq-x$ implique, pour $0\leq k<n$,
\[
0\leq\left(1-\frac{k}{n}\right)^{\alpha n}\leq\e^{-\alpha k}.
\]
Par conséquent,
\[
\sum_{k>K_n}^{n-1}\left(1-\frac{k}{n}\right)^{\alpha n}
\leq\frac{\e^{-\alpha(K_n+1)}}{1-\e^{-\alpha}}\to0,
\]
et la même majoration vaut pour $\sum_{k>K_n}\e^{-\alpha k}$. Il s'ensuit
\[
A_n-\sum_{k=0}^{\infty}\e^{-\alpha k}\to0.
\]
La somme géométrique vaut $(1-\e^{-\alpha})^{-1}$, ce qui conclut.
\end{proofbox}

\begin{remarkbox}
Tout $\lambda\in(0,1)$ convient grâce au facteur décroissant
$\e^{-\alpha k}$. Une estimation uniforme plus grossière, sommée terme à
terme sans conserver ce facteur, imposerait artificiellement une restriction
sur $\lambda$.
\end{remarkbox}

\subsection{Deuxième preuve : convergence dominée sur $\N$}

\begin{proofbox}
Prolongeons les termes par zéro et posons, pour $k\in\N$,
\[
f_n(k):=\mathbf 1_{\{k<n\}}\left(1-\frac{k}{n}\right)^{\alpha n}.
\]
Pour tout $k$ fixé, $f_n(k)\to\e^{-\alpha k}$. En outre,
$0\leq f_n(k)\leq\e^{-\alpha k}$, et
$\sum_{k\geq0}\e^{-\alpha k}<\infty$. Le théorème de convergence dominée
sur l'espace mesuré $(\N,\mathcal P(\N),\text{comptage})$ entraîne
\[
A_n=\sum_{k=0}^{\infty}f_n(k)\longrightarrow
\sum_{k=0}^{\infty}\e^{-\alpha k}=\frac{1}{1-\e^{-\alpha}}.
\]
\end{proofbox}

\subsection{Troisième preuve : encadrement par deux séries géométriques}

\begin{proofbox}
Fixons $0<\eta<1$. Pour $0\leq k\leq\eta n$, l'inégalité
\[
-\frac{x}{1-x}\leq\log(1-x)\leq-x\qquad(0\leq x<1)
\]
donne
\[
\e^{-\alpha k/(1-\eta)}
\leq\left(1-\frac{k}{n}\right)^{\alpha n}leq\e^{-\alpha k}.
\]
Ainsi
\[
\sum_{k=0}^{\lfloor\eta n\rfloor}\e^{-\alpha k/(1-\eta)}
\leq A_n\leq\sum_{k=0}^{n-1}\e^{-\alpha k}.
\]
En faisant $n\to\infty$, on obtient
\[
\frac1{1-\e^{-\alpha/(1-\eta)}}
\leq\liminf A_n\leq\limsup A_n\leq\frac1{1-\e^{-\alpha}}.
\]
Enfin $\eta\downarrow0$ fait tendre le membre de gauche vers le membre de
droite; la limite de $A_n$ existe et possède la valeur annoncée.
\end{proofbox}

\section{Exercice 26 : puissances de factorielles}

\begin{theorembox}
Soit $\alpha>0$. Lorsque $n\to\infty$,
\[
 S_n:=\sum_{p=1}^{n}(p!)^{-\alpha/n}sim
 \frac1\alpha\frac{n}{\log n}.
\]
\end{theorembox}

Posons $N_n=n/\log n$. Les inégalités intégrales
\[
p\log p-p+1\leq\log(p!)\leq p\log p
\]
entraînent
\begin{equation}
p^{-\alpha p/n}\leq(p!)^{-\alpha/n}
\leq \e^{\alpha(p-1)/n}p^{-\alpha p/n}.\tag{1}
\end{equation}

\subsection{Première preuve : somme modèle et coupure $\lfloor\eps n\rfloor$}

\begin{lemmabox}[title=Somme modèle]
Si $T_n=\sum_{p=1}^{n}p^{-\alpha p/n}$, alors
\[
\frac{T_n}{N_n}\longrightarrow\frac1\alpha.
\]
\end{lemmabox}

\begin{proofbox}
Fixons $0<\delta<M$. Pour $p=u_pN_n$ avec
$u_p\in[\delta,M]$,
\[
\frac{\alpha p\log p}{n}
=\alpha u_p\left(1-\frac{\log\log n}{\log n}
+\frac{\log u_p}{\log n}\right)=\alpha u_p+o(1),
\]
uniformément sur cet intervalle. Les sommes de Riemann donnent donc
\[
\frac1{N_n}\sum_{\delta N_n\leq p\leq MN_n}p^{-\alpha p/n}
\to\int_\delta^M\e^{-\alpha u}\,du.
\]
La partie $p<\delta N_n$ contribue au plus $\delta+o(1)$ après division par
$N_n$. Pour $p\geq MN_n$, et $n$ assez grand, on traite d'abord l'intervalle
$[MN_n,\sqrt n]$, qui est vide puisque $N_n/\sqrt n\to\infty$. Puis, pour
$p\geq\sqrt n$, $\log p\geq\frac12\log n$, donc
\[
p^{-\alpha p/n}\leq\e^{-\alpha p/(2N_n)}.
\]
La somme géométrique correspondante, divisée par $N_n$, est
$O_\alpha(\e^{-\alpha M/2})$. En faisant successivement
$n\to\infty$, $\delta\downarrow0$ et $M\to\infty$, on obtient
$T_n/N_n\to\int_0^\infty\e^{-\alpha u}\,du=1/\alpha$.
\end{proofbox}

\begin{proofbox}[title=Comparaison avec la factorielle]
Fixons $\eps\in(0,1)$. Pour $p\leq\eps n$, (1) fournit
\[
p^{-\alpha p/n}\leq(p!)^{-\alpha/n}leq
\e^{\alpha\eps}p^{-\alpha p/n}.
\]
Pour $p>\eps n$, l'inégalité inférieure sur $\log(p!)$ montre, pour $n$
assez grand,
\[
(p!)^{-\alpha/n}\leq\exp\left(-\frac{\alpha p\log n}{2n}\right)
=\e^{-\alpha p/(2N_n)}.
\]
La somme de cette queue est $o(N_n)$; il en va de même pour la queue de
$T_n$. Le lemme entraîne alors
\[
\frac1\alpha\leq\liminf\frac{S_n}{N_n}leq
\limsup\frac{S_n}{N_n}\leq\frac{\e^{\alpha\eps}}\alpha.
\]
Le passage $\eps\downarrow0$ conclut.
\end{proofbox}

\subsection{Deuxième preuve : approximation uniforme sur l'échelle naturelle}

\begin{proofbox}
Soient $0<\delta<M$. Pour $p=uN_n$, $u\in[\delta,M]$, la formule de
Stirling uniforme sous la forme $\log(p!)=p\log p-p+O(\log p)$ donne
\[
\frac{\alpha}{n}\log(p!)=\alpha u+o(1)
\]
uniformément en $u\in[\delta,M]$. Par conséquent,
\[
\frac1{N_n}\sum_{\delta N_n\leq p\leq MN_n}(p!)^{-\alpha/n}
\to\int_\delta^M\e^{-\alpha u}\,du.
\]
La zone $p<\delta N_n$ est majorée par $\delta N_n$. La queue $p>MN_n$
est contrôlée exactement comme dans la première preuve par
$\e^{-\alpha p/(2N_n)}$, dès que $p\geq\sqrt n$; l'éventuel intervalle
$[MN_n,\sqrt n]$ est vide pour $n$ assez grand. Après division par $N_n$,
sa contribution est $O_\alpha(\e^{-\alpha M/2})$. Les passages
$\delta\downarrow0$ et $M\to\infty$ donnent
\[
\frac{S_n}{N_n}\to\int_0^\infty\e^{-\alpha u}\,du=\frac1\alpha.
\]
\end{proofbox}

\subsection{Troisième preuve : comparaison intégrale monotone}

\begin{proofbox}
La fonction $x\mapsto\log\Gamma(x+1)$ est croissante sur $[1,\infty)$ à
partir d'un point fixe; modifier un nombre fini de termes produit une erreur
$O(1)=o(N_n)$. Ainsi, par encadrement somme--intégrale,
\[
S_n=\int_1^n\exp\left(-\frac{\alpha}{n}\log\Gamma(x+1)\right)dx+o(N_n).
\]
Effectuons $x=uN_n$. Pour $u$ dans tout compact de $(0,\infty)$, Stirling
donne uniformément
\[
\frac{\alpha}{n}\log\Gamma(uN_n+1)=\alpha u+o(1).
\]
L'intégrale normalisée converge donc sur $[\delta,M]$ vers
$\int_\delta^M\e^{-\alpha u}du$. La partie $u<\delta$ est au plus $\delta$.
Pour $u>M$, l'estimation
$\log\Gamma(x+1)\geq\frac12x\log n$ valable lorsque
$x\geq\sqrt n$ donne un majorant intégrable $\e^{-\alpha u/2}$; or
$MN_n>\sqrt n$ pour $n$ assez grand. La queue normalisée est donc au plus
$\int_M^\infty\e^{-\alpha u/2}du$. On conclut en faisant
$n\to\infty$, puis $\delta\downarrow0$ et $M\to\infty$.
\end{proofbox}

\section{Conclusion comparative}

\begin{propositionbox}[title=Choix des méthodes]
La preuve par $n^\lambda$ de l'exercice 25 et la preuve par
$\lfloor\eps n\rfloor$ de l'exercice 26 suivent directement les indications
de Dieudonné et gardent un contrôle quantitatif des queues. Les preuves par
convergence dominée et par Stirling sont plus courtes, mais mobilisent des
théorèmes plus généraux. Les encadrements géométriques et intégraux offrent
des alternatives élémentaires et conceptuellement distinctes.
\end{propositionbox}

\end{document}
